\documentclass[11pt,a4paper]{article}
\usepackage[margin=24mm]{geometry}
\usepackage{amsmath,amssymb,amsthm}
\usepackage[T1]{fontenc}
\usepackage{lmodern}
\usepackage{hyperref}
\newtheorem{theorem}{Theorem}[section]
\newtheorem{proposition}[theorem]{Proposition}
\newtheorem{corollary}[theorem]{Corollary}
\title{Research as Production Context\\[4pt]\large Conditional labor requirements, reusable capabilities,\\and normalized occupational conversion}
\author{}
\date{24 October 2026}
\begin{document}
\maketitle
\begin{abstract}
A production recipe consumes inputs and labor but may also depend on reusable knowledge. We describe these two kinds of requirement by a nonnegative matrix inverse and a Boolean path closure. They determine occupational labor, supporting capabilities, and the fresh cost of projects. The same constructions give inventory identities, boundary reductions, sensitivity formulas, and error bounds. An optional equal-residual convention determines normalized occupational weights through a second nonnegative inverse. We prove the results and give small examples that identify their limits. All conclusions are conditional on the selected recipes, units, and stated assumptions.
\end{abstract}

\section{The model and its solution}
Fix $n\geq1$ outputs with specified units and selected recipes, and $m\geq1$ occupations. Let $A\geq0$ be the $n\times n$ input matrix and $H\geq0$ the $m\times n$ direct-labor matrix: $A_{ij}$ is input $i$ per unit of output $j$, and $H_{rj}$ is category-$r$ labor per unit of output $j$. Vectors are columns; inequalities are entrywise. Write $w\gg0$ for strictly positive components, $Aw\ll w$ for componentwise strict inequality, and $\mathbf1$ for an all-ones vector. The spectral radius $\rho(A)$ includes complex eigenvalues; $\|\cdot\|_\infty$ is the maximum absolute row-sum norm.

These data describe one production context. Recipe eligibility, capacity, scheduling, and project success remain separate conditions. Reusable capabilities are recorded by a graph in Section 2; creating them is a project whose costs are accounted for explicitly. No historical cost allocation or occupational weighting is implicit in $H$.

\begin{theorem}[Productivity]
For $A\geq0$, the following are equivalent: (i) $\rho(A)<1$; (ii) $\sum_{k\geq0}A^k$ converges; (iii) $I-A$ has a nonnegative inverse; (iv) some $w\gg0$ satisfies $Aw\ll w$; (v) every $y\geq0$ has a unique nonnegative solution of $x=Ax+y$. Call $A$ productive when these hold. Then
\begin{equation}
N=(I-A)^{-1}=\sum_{k\geq0}A^k\geq I,\qquad x=Ny.
\end{equation}
\end{theorem}
\begin{proof}
For (i)$\Rightarrow$(ii), choose $\rho(A)<r<1$. The spectral-radius formula gives $\|A^k\|\leq r^k$ eventually, so the series converges. For (ii)$\Rightarrow$(iii), take limits in $(I-A)\sum_{k=0}^{K}A^k=I-A^{K+1}=\sum_{k=0}^{K}A^k(I-A)$. Its nonnegative sum is a two-sided inverse and is at least $I$.

For (iii)$\Rightarrow$(iv), put $w=N\mathbf1$. Then $w=Aw+\mathbf1\geq\mathbf1$. For (iv)$\Rightarrow$(i), put $P=\operatorname{diag}(w)$. Similarity and the row-sum bound give
\begin{equation}
\rho(A)=\rho(P^{-1}AP)\leq\|P^{-1}AP\|_\infty=\max_i\frac{(Aw)_i}{w_i}<1.
\end{equation}
Finally, (iii) gives the unique solution $x=Ny\geq0$; conversely, a nonnegative solution for $y=\mathbf1$ is a witness for (iv).
\end{proof}

\begin{corollary}[Labor, final demand, and projects]
For productive $A$, the unique solution of $G=H+GA$ is
\begin{equation}
G=HN=\sum_{k\geq0}HA^k\geq0.
\end{equation}
Final demand $y\geq0$ requires $Gy=HNy$ occupational hours. A project with direct labor $h\geq0$ and freshly reproduced inputs $a\geq0$ requires $\ell=h+Ga$. For specified weights $\omega\gg0$, $v=G^{\mathsf T}\omega$ is the unique solution of $v=H^{\mathsf T}\omega+A^{\mathsf T}v$, and the scalar accounts are $v^{\mathsf T}y=\omega^{\mathsf T}Gy$ and $\omega^{\mathsf T}\ell$.
\end{corollary}
\begin{proof}
Multiply $G(I-A)=H$ by $N$; the series follows by continuity. Reproducing $a$ needs gross output $Na$, hence labor $HNa$, in addition to $h$. Transposing the labor equation gives the scalar equation, whose inverse is $N^{\mathsf T}$. Taking column $j$ also gives the unit requirement $Ge_j=He_j+GAe_j$.
\end{proof}

A balance for one particular $H$ does not establish productivity: $A=[2]$, $H=[0]$ admits $G=[0]$. Nor does the balance determine $\omega$. Section 4 supplies an optional weighting convention.

\section{Capabilities and development projects}
Let $B$ be a Boolean $n\times n$ dependency matrix: $B_{ij}=1$ means that output $j$ depends on upstream output $i$. Let $D$ be a Boolean $q\times n$ matrix of direct capability labels, allowing $q=0$, and require $\operatorname{supp}(A)\leq B$, where support marks positive entries. Boolean multiplication $\odot$ uses AND along a path and OR across paths. All other products use real arithmetic.

\begin{theorem}[Path closure and numerical support]
With Boolean powers, define
\begin{equation}
B^*=\bigvee_{k=0}^{n-1}B^k,\qquad R=D\odot B^*.
\end{equation}
Then $R$ is the least solution of $R=D\vee(R\odot B)$. Iterating from $R_0=D$ reaches it after at most $n-1$ rounds. If $A$ is productive and $\widehat D$ is the real zero-one version of $D$, then
\begin{equation}
\operatorname{supp}(N)=\operatorname{supp}(A)^*,\qquad
\operatorname{supp}(\widehat D N)=D\odot\operatorname{supp}(A)^*\leq R.
\end{equation}
The last inequality is an equality when $B=\operatorname{supp}(A)$.
\end{theorem}
\begin{proof}
Deleting cycles reduces a path to at most $n-1$ edges, with length zero covering each vertex itself. Splitting off or appending a last edge proves the fixed-point equation. Any fixed point contains $D$ and, inductively, every $D\odot B^k$, proving leastness. The $s$th iterate contains exactly the labeled paths of length at most $s$.

For nonnegative matrices, a product entry is positive exactly when one summand is positive. Thus $(A^k)_{ij}>0$ records a length-$k$ support path. An entry of the nonnegative series for $N$ is positive exactly when some term is positive. This proves both support identities; adding declared edges can only enlarge closure.
\end{proof}

Leastness is essential: for $D=0$, $B=I$, every Boolean matrix is a fixed point. Also, with $A=0$, $D=(1,0)$, and a declared edge $1\to2$, numerical support is $(1,0)$ but $R=(1,1)$. Closure records upstream support, while local eligibility still asks whether a recipe's direct labels are accessible to its producer.

\begin{proposition}[Missing projects and terminal outputs]
For $r\geq1$ projects, let $Q_{ij}=1$ mean that project $j$ needs result $i$. Let $b$ mark missing results and $\tau$ mark targets. The iteration
\begin{equation}
u_0=b\wedge\tau,\qquad u_{s+1}=b\wedge\bigl(\tau\vee(Q\odot u_s)\bigr)
\end{equation}
reaches its least fixed point within $r-1$ rounds. It marks the missing targets and prerequisites reached along paths entirely within the missing set.

For productive $A$, appending an output used by no selected recipe gives
\begin{equation}
\overline A=\begin{pmatrix}A&a\\0&0\end{pmatrix},\quad
(I-\overline A)^{-1}=\begin{pmatrix}N&Na\\0&1\end{pmatrix},\quad
\overline G=\begin{pmatrix}G&h+Ga\end{pmatrix},
\end{equation}
where $a,h\geq0$ and $\overline H=(H\ h)$. Hence $\overline A$ is productive. With $G$ fixed, $a_\varepsilon,h_\varepsilon\to0$ implies $h_\varepsilon+Ga_\varepsilon\to0$.
\end{proposition}
\begin{proof}
Apply the path argument to backward prerequisite paths, restricting every vertex to the missing set. The mask stops expansion at satisfied results. For the terminal output, multiply the displayed block matrices to verify the nonnegative inverse, then multiply by $\overline H$. The limit follows by continuity.
\end{proof}

A successful project adds its outcome capabilities; a failed attempt may still incur costs. For a fixed context with fresh inputs, attempting each selected project once sums costs over distinct projects. Two projects costing 6 and 8 with a shared prerequisite costing 4 therefore require 18, or 14 if the prerequisite is already satisfied. Closure alone does not make a cycle of missing prerequisites executable. A copying requirement may tend to zero while its release capability remains necessary.

\section{Inventories}
At stage $t$, write $x_t$ for gross production, $y_t$ for deliveries, $s_t$ for opening stocks, and $z_t$ for project attempts. Let $J_t$ contain their ordinary inputs and $K_t$ their direct labor. Assume $s_{t+1}=s_t+(I-A_t)x_t-y_t-J_tz_t$ and productivity of each $A_t$.

\begin{proposition}[Chronological reconciliation]
For $G_t=H_t(I-A_t)^{-1}$, actual stage labor is
\begin{equation}
H_tx_t+K_tz_t=G_t(y_t+J_tz_t+s_{t+1}-s_t)+K_tz_t.
\end{equation}
If $A,H$ are constant, total labor over $T$ stages is
\begin{equation}
L=\sum_{t=0}^{T-1}\bigl[Gy_t+(K_t+GJ_t)z_t\bigr]+G(s_T-s_0).
\end{equation}
\end{proposition}
\begin{proof}
Rearrange the stock balance and use $G_t(I-A_t)=H_t$. Sum; when $G$ is fixed, its stock increments telescope.
\end{proof}

An input produced in one hour and later consumed by a project using two direct hours contributes three hours in total. Adding both its production hour and the project's three-hour fresh requirement would count that input twice. The stock correction prevents this. With changing $G_t$, retain each weighted increment: $G_0=1$, $G_1=2$, and stocks $(0,1,0)$ give $1-2=-1$, despite zero net stock change. Signed stock changes and real attempt quantities are permitted by these identities.

\section{Optional occupational weights}
For a productive context, supply provisioning goods $S\geq0$, training goods $E\geq0$, direct teaching labor $U\geq0$, and positive finite lifetimes $\Gamma=\operatorname{diag}(\gamma_i)$. The first two matrices are $n\times m$; $U,\Gamma$ are $m\times m$. Define
\begin{equation}
F=S+E\Gamma^{-1},\qquad T=U\Gamma^{-1},\qquad M=GF+T.
\end{equation}
These data specify reproduction burdens per productive time unit, with goods and direct teaching counted separately. Fix reference occupation $b$. The equal-residual convention requires $\alpha\gg0$, $\kappa>0$, and
\begin{equation}
(I-M^{\mathsf T})\alpha=\kappa\mathbf1,\qquad\alpha_b=1.
\end{equation}

\begin{theorem}[Normalized weights]
Such a pair exists exactly when $\rho(M)<1$, and then it is unique:
\begin{equation}
W=(I-M^{\mathsf T})^{-1},\quad\beta=W\mathbf1\geq\mathbf1,\quad
\alpha=\frac{\beta}{\beta_b},\quad\kappa=\frac1{\beta_b}.
\end{equation}
The ratios $r_{ij}=\alpha_i/\alpha_j=\beta_i/\beta_j$ are independent of the reference. Changing reference to $c$ divides $\alpha$ and $\kappa$ by $\alpha_c$. Equal column sums $M^{\mathsf T}\mathbf1=s\mathbf1$, with $s<1$, give $\alpha=\mathbf1$ and $\kappa=1-s$.
\end{theorem}
\begin{proof}
Productivity of $M^{\mathsf T}$ gives $W\geq I$. Every solution has $\alpha=\kappa\beta$; normalization forces the displayed pair. Conversely, a positive pair gives $M^{\mathsf T}\alpha\ll\alpha$, proving productivity. Ratios cancel the normalizing denominator. For equal column sums, $\beta=(1-s)^{-1}\mathbf1$ solves the equation uniquely.
\end{proof}

This is an additional convention, not an identification of weights with wages. Ordinary productivity is insufficient: $A=0$, $H=S=1$, $E=U=0$ gives $M=1$ and no positive residual.

\begin{proposition}[The augmented system]
For productive $A$, let $\mathcal A=\left(\begin{smallmatrix}A&F\\H&T\end{smallmatrix}\right)$. Then $\mathcal A$ is productive exactly when $M$ is. When these are productive, write $Z=(I-M)^{-1}$; the inverse is
\begin{equation}
(I-\mathcal A)^{-1}=\begin{pmatrix}N+NFZHN&NFZ\\ZHN&Z\end{pmatrix}.
\end{equation}
The unique normalized scalar account is $v=G^{\mathsf T}\alpha$, equivalently
\begin{equation}
\begin{pmatrix}v\\\alpha\end{pmatrix}=\mathcal A^{\mathsf T}\begin{pmatrix}v\\\alpha\end{pmatrix}+\begin{pmatrix}0\\\kappa\mathbf1\end{pmatrix},\qquad\alpha_b=1,
\end{equation}
with $\alpha\gg0$, $\kappa>0$. Projects have scalar requirement $\alpha^{\mathsf T}(h+Ga)$.
\end{proposition}
\begin{proof}
For arbitrary right-hand sides $p,q$, eliminate $u$ from $(I-A)u-Fz=p$ and $-Hu+(I-T)z=q$. This gives $z=Z(q+HNp)$ and $u=N(p+Fz)$, proving the inverse formula and its nonnegativity when $M$ is productive. Conversely, a positive witness $(u,z)$ for $\mathcal A$ gives $NFz\leq u$ and hence $Mz\leq Hu+Tz\ll z$. Finally, eliminating $v$ from the transposed equations gives exactly the equal-residual equation.
\end{proof}

\section{Boundaries and units}
Partition outputs into nonempty retained nodes $L$ and external nodes $O$, with corresponding blocks of $A$, columns of $H,G,D$, and rows of $F$. All formulas below retain the input-per-output orientation: $A_{LO}$ means retained inputs used by external production. Empty-block cases simply omit the corresponding terms.

\begin{theorem}[Exact boundary reduction]
If $A$ is productive, so are $A_{OO}$ and $A_{\mathrm{eff}}$ below. Set
\begin{equation}
N_O=(I-A_{OO})^{-1},\quad
A_{\mathrm{eff}}=A_{LL}+A_{LO}N_OA_{OL},\quad
H_{\mathrm{eff}}=H_L+H_ON_OA_{OL}.
\end{equation}
Then
\begin{equation}
G_L=H_{\mathrm{eff}}(I-A_{\mathrm{eff}})^{-1},\qquad
G_O=(H_O+G_LA_{LO})N_O.
\end{equation}
For final demand $(y_L,y_O)$, $x_L=(I-A_{\mathrm{eff}})^{-1}(y_L+A_{LO}N_Oy_O)$ and total labor is $H_{\mathrm{eff}}x_L+H_ON_Oy_O$. Conversion is preserved by
\begin{equation}
F_{\mathrm{eff}}=F_L+A_{LO}N_OF_O,\quad T_{\mathrm{eff}}=T+H_ON_OF_O,\quad
G_LF_{\mathrm{eff}}+T_{\mathrm{eff}}=M.
\end{equation}
Capability closure is preserved by the Boolean formulas
\begin{equation}
\begin{aligned}
B_{\mathrm{eff}}&=B_{LL}\vee(B_{LO}\odot B_{OO}^*\odot B_{OL}),\\
D_{\mathrm{eff}}&=D_L\vee(D_O\odot B_{OO}^*\odot B_{OL}),\qquad
R_L=D_{\mathrm{eff}}\odot B_{\mathrm{eff}}^*.
\end{aligned}
\end{equation}
\end{theorem}
\begin{proof}
A witness $(w_L,w_O)$ for $A$ gives $A_{OO}w_O\ll w_O$ and $N_OA_{OL}w_L\leq w_O$. Therefore $A_{\mathrm{eff}}w_L\leq A_{LL}w_L+A_{LO}w_O\ll w_L$. Solve the external labor equation for $G_O$ and substitute into the retained equation. The production identity $x_O=N_O(A_{OL}x_L+y_O)$ gives the demand formulas. Substitution gives $G_LF_{\mathrm{eff}}+T_{\mathrm{eff}}=G_LF_L+G_OF_O+T=M$.

For capabilities, compress each external excursion of a path into an effective edge, and any initial external labeled segment into an effective label. Expanding those segments reverses the construction. Thus full and reduced graphs have exactly the same labeled reachability on $L$.
\end{proof}

Consequently the reduced system preserves normalized weights, ratios, and retained scalar values whenever conversion exists. Effective coefficients include eliminated labor and inputs; simply deleting external rows, or external final demand, is not this reduction.

\begin{proposition}[Changes of units]
Assume $A$ productive, and $M$ productive when weights are used. For positive diagonal physical-unit matrix $P$, set $A'=P^{-1}AP$, $H'=HP$, $y'=P^{-1}y$, $S'=P^{-1}S$, and $E'=P^{-1}E$, leaving $U,\Gamma$ fixed. Then
\begin{equation}
N'=P^{-1}NP,\quad G'=GP,\quad G'y'=Gy,\quad M'=M,\quad v'=Pv.
\end{equation}
Productivity, support, conversion weights, and scalar totals are preserved.

If one new time unit in occupation $i$ is $t_i>0$ old units, put $J=\operatorname{diag}(t_i)$ and keep physical units fixed. Use
\begin{equation}
H'=J^{-1}H,\quad S'=SJ,\quad E'=E,\quad U'=J^{-1}U,\quad\Gamma'=J^{-1}\Gamma.
\end{equation}
Then $G'=J^{-1}G$, $F'=FJ$, $T'=J^{-1}TJ$, $M'=J^{-1}MJ$, and the covariant weights $\alpha'=J\alpha$ satisfy
\begin{equation}
\alpha'-(M')^{\mathsf T}\alpha'=\kappa J\mathbf1,\qquad
(G')^{\mathsf T}\alpha'=G^{\mathsf T}\alpha.
\end{equation}
Rebasing at $b$ divides $\alpha'$ and $\kappa$ by $t_b$ and divides scalar values by $t_b$, keeping source $J\mathbf1$. Ratios per new unit are $(t_i/t_j)r_{ij}$.
\end{proposition}
\begin{proof}
Substitute $I-A'=P^{-1}(I-A)P$ for physical units. For time units use $(\Gamma')^{-1}=\Gamma^{-1}J$ and $(M')^{\mathsf T}=JM^{\mathsf T}J^{-1}$. The formulas follow by cancellation; positive diagonal factors preserve support.
\end{proof}

The residual source must change with occupational time units. Keeping a numerical all-ones source after unequal rescaling would impose a different weighting convention.

\section{Changes and error bounds}
\begin{proposition}[Finite changes and derivatives]
For two productive systems, write $N'=(I-A')^{-1}$. Then
\begin{equation}
G'-G=\bigl(H'-H+G(A'-A)\bigr)N'.
\end{equation}
If $0\leq A'\leq A$ and $0\leq H'\leq H$, productivity of $A$ implies that of $A'$ and $G'\leq G$. For differentiable coefficients at a productive point, in fixed coordinates,
\begin{equation}
dN=N\,dA\,N,\qquad dG=(dH+G\,dA)N,\qquad
\frac{\partial G_{rj}}{\partial A_{ik}}=G_{ri}N_{kj}.
\end{equation}
When conversion is productive, also
\begin{equation}
\begin{aligned}
dM&=dG\,F+G\,dF+dT,\qquad d\beta=W(dM)^{\mathsf T}\beta,\\
d\alpha&=\frac{d\beta-\alpha\,d\beta_b}{\beta_b},\qquad dv=(dG)^{\mathsf T}\alpha+G^{\mathsf T}d\alpha.
\end{aligned}
\end{equation}
\end{proposition}
\begin{proof}
Subtract the balance equations and multiply by $N'$ to obtain the finite identity. A witness for $A$ is a witness for $A'\leq A$; the identity then gives $G'\leq G$. The inverse map is differentiable on the open set of invertible matrices. Differentiating $(I-A)N=I$ gives $dN$, and the product rule gives $dG$; setting $dA=e_i e_k^{\mathsf T}$ gives the coefficient formula with $H$ fixed. Differentiate $(I-M^{\mathsf T})\beta=\mathbf1$, then apply the product and quotient rules, using $\beta_b>0$.
\end{proof}

These derivatives concern the inverse formulas, not changes of catalogue or feasibility regime. Ordered productive matrices $M\leq M'$ give $\beta\leq\beta'$ by the nonnegative series, but normalization can reverse an inequality: $M=0$ and $M'=\operatorname{diag}(1/2,0)$ give reference-1 weights $(1,1)$ and $(1,1/2)$.

\begin{proposition}[Certified approximations]
Assume $A$ productive, and $M$ productive for the conversion bounds. Let $P=\operatorname{diag}(w)$ for a positive witness of $A$, and put $\eta=\|P^{-1}AP\|_\infty<1$. With $G^{[k]}=H\sum_{i=0}^{k}A^i$, one has $0\leq G^{[k]}\leq G$ and
\begin{equation}
\|(G-G^{[k]})P\|_\infty\leq\|HP\|_\infty\frac{\eta^{k+1}}{1-\eta}.
\end{equation}
For an approximation $\widetilde\beta$, let $r=\mathbf1-(I-M^{\mathsf T})\widetilde\beta$. Then
\begin{equation}
\beta-\widetilde\beta=Wr,\qquad |\beta-\widetilde\beta|\leq W|r|.
\end{equation}
Bounds $0<l_i\leq\beta_i\leq u_i$ imply $l_i/u_j\leq r_{ij}\leq u_i/l_j$.
\end{proposition}
\begin{proof}
After scaling, the labor error is the tail $HP\sum_{i=k+1}^{\infty}(P^{-1}AP)^i$. Bound its norm by the scalar geometric tail. Subtracting the approximate conversion equation from the exact one gives $Wr$; use $W\geq0$ and the triangle inequality. The ratio bounds follow by division by positive bounds.
\end{proof}

Thus a small residual needs the factor $W$ to certify accuracy. An unscaled row-sum test is only sufficient for productivity: a strictly upper triangular matrix is nilpotent regardless of how large its entries are.

\section{Small examples}
For two outputs and two occupations, take
\begin{equation}
A=\begin{pmatrix}0&2\\0&0\end{pmatrix},\quad
H=\begin{pmatrix}3&1\\0&2\end{pmatrix},\quad
G=H(I+A)=\begin{pmatrix}3&7\\0&2\end{pmatrix}.
\end{equation}
The second output requires occupational hours $(7,2)^{\mathsf T}$; supplied weights $(1,3)^{\mathsf T}$ assign it scalar requirement 13.

For a diamond with unit coefficients from node 1 to nodes 2 and 3, and from both to node 4, let only node 1 use one direct hour. Material requirements are $(1,1,1,2)$ while a capability labeling node 1 propagates to $(1,1,1,1)$. Material paths add; reusable labels combine by OR. Similarly, development costing 1000 hours followed by 1000 copies at $1/100$ hour each totals 1010 hours, while each conditional copy requirement remains $1/100$.

For conversion, take
\begin{equation}
M=\begin{pmatrix}1/5&3/10\\1/20&1/4\end{pmatrix},\qquad
W=\begin{pmatrix}50/39&10/117\\20/39&160/117\end{pmatrix}.
\end{equation}
Direct multiplication verifies $W=(I-M^{\mathsf T})^{-1}\geq0$, giving
\begin{equation}
\beta=\frac1{117}\begin{pmatrix}160\\220\end{pmatrix},\qquad
\alpha=\begin{pmatrix}1\\11/8\end{pmatrix},\qquad\kappa=\frac{117}{160}.
\end{equation}
A project requiring $(15,22)^{\mathsf T}$ hours has scalar requirement $15+(11/8)22=181/4$. Changing the second time unit to eight hours gives weights $(1,11)^{\mathsf T}$ and residual source $(1,8)^{\mathsf T}$, preserving scalar requirements.

The weighting assumption matters. Replacing the residual source by $c=(1,2)^{\mathsf T}$ gives $Wc=(170,380)^{\mathsf T}/117$, hence reference-1 weight $38/17$ and residual $117/170$. Omitting goods burdens and using only $T=\left(\begin{smallmatrix}0&0\\3/100&1/10\end{smallmatrix}\right)$ would instead give weight $100/93$. Equal clock hours alone also need not receive equal weights: $A=0$, $H=I$, $E=U=0$, and $S=\operatorname{diag}(1/5,2/5)$ give $\alpha=(1,4/3)^{\mathsf T}$.

\section{Verification note}
This exposition consolidates the results of the supplied Research as Production Context, audited edition dated 21 October 2026. The accompanying Lean 4.35.0 and Mathlib v4.35.0 development contains 191 named theorems and lemmas, with 477 declarations covered by its axiom audit and no proof placeholders. The source project is the formal certificate; the README maps the results. Empirical validity and the choice of weighting convention remain separate from the proved identities.
\end{document}
